\section{Experimental design}
\subsection{Systems and operating points}
The dense reference is AudioLDM-M-Full~\cite{liu2023audioldm}, with $415.96$\,M U-Net parameters. We study the two released pruning severities from Singh et al.~\cite{singh2026pruning}. Severity~1 retains $145.67$\,M parameters after $65.0\%$ pruning, while severity~2 retains $71.08$\,M after $82.9\%$ pruning. At each severity, P denotes the checkpoint before recovery and \PFT{} the released checkpoint after $10^{6}$ AudioCaps fine-tuning steps. Their paired difference measures what the released recovery stage adds within the pruned architecture. It does not by itself identify an effect caused specifically by pruning.

The primary duration comparison uses $3.84$ and $10.24$\,s. Severity~2 adds $5.12$ and $7.68$\,s, producing a four-point sweep. AudioCaps \textsc{test} supplies the in-domain prompts~\cite{kim2019audiocaps}. Clotho provides a held-out sound-event domain with a different caption distribution~\cite{drossos2020clotho}. A second held-out battery uses hip-hop and rap captions from MusicCaps~\cite{agostinelli2023musiclm}. The main severity~2 AudioCaps analysis uses $192$ prompts, Clotho uses $96$, and the extended hip-hop battery uses $127$.

Generation uses DDIM~\cite{song2021ddim} with $50$ steps, guidance $2.5$ and $\eta=0$. A pre-specified robustness check repeats the severity~2 endpoint comparison with the published DDIM $200$, guidance $3.5$ recipe.

\subsection{Interventions and weight conventions}
The mechanism and boundary checks were frozen before their own generation and scoring. They test explanations of the primary finding but are not part of its original confirmatory family.

The central intervention is symmetric. Starting from the same severity~2 P checkpoint, we fine-tune the full U-Net for $20{,}000$ optimizer updates at either $3.84$ or $10.24$\,s while holding every other training choice fixed. Both runs use AdamW with a constant learning rate of $10^{-4}$, default $\beta$ values $(0.9,0.999)$, weight decay $0.01$, no scheduler and effective batch size $2$. The two resulting checkpoints are evaluated at both durations on the same $192$ prompts with common generation noise. If recovery specializes to its training duration, moving training from $3.84$ to $10.24$\,s should increase the long-minus-short interaction, so
\begin{equation}
\Delta J = J_{10.24}-J_{3.84}>0.
\end{equation}

A reduced dense $2\times2$ repeats this exact $20{,}000$-step recipe from AudioLDM-M-Full. Because the experimental exports are raw weights, we also evaluate the raw released dense checkpoint on the same prompts. Raw and EMA baselines differ by at most about $0.02$ CLAP at either duration, which rules out weight convention as an explanation for the much larger changes after fine-tuning. This dense experiment probes the short adaptation recipe only. It cannot replace the unavailable dense checkpoint after the released $10^{6}$-step recovery, and no lower-learning-rate dense run was performed.

The severity~2 P used for the symmetric intervention is derived from dense EMA weights, whereas both $20{,}000$-step fine-tuned checkpoints are raw exports. This convention mismatch cancels exactly in $\Delta J$. On the dense model at step zero, its duration-dependent offset is $-0.019$ \ci{-0.050}{+0.013}, so we treat the individual $J$ values with this caveat and base the specialization test on their paired difference.

Further checks extend the released recovery to $15.36$\,s on $96$ matched prompts, enlarge severity~1 with a disjoint $96$-prompt sample, split the severity~2 AudioCaps draw into two outcome-blind halves, and complete dense hip-hop anchors for all $127$ prompts. These analyses probe response shape, prompt-sample stability and metric-floor interpretations.

\subsection{Pairing, endpoints and estimands}
Within each operating point, compared checkpoints receive identical diffusion noise $x_T$ for a given prompt. The prompt is the statistical unit, and $95\%$ intervals use a prompt-level percentile bootstrap with $B=10^{4}$ resamples. This common-random-number design makes $R$ a paired effect rather than a difference between independently sampled generation sets.

The primary endpoint is fused CLAP cosine from \texttt{laion/\allowbreak clap-htsat-fused}, revision \texttt{365dea6e}~\cite{wu2023clap}. Human-CLAP provides a second alignment scorer. KL is measured against matched real references, with positive recovery meaning that \PFT{} reduces KL relative to P. For PANNs, \emph{top-10 capture} is the number of ground-truth AudioSet labels from the reference clip that appear among the generated clip's ten highest-scoring PANNs classes. PANNs capture recovery is the paired increase in that count from P to \PFT{}.

For operating point $o$, recovery gain is
\begin{equation}
R(o)=\mathbb{E}_{i}\left[S(\PFT,i,o)-S(P,i,o)\right].
\end{equation}
The duration interaction is $J=R(10.24\,\mathrm{s})-R(3.84\,\mathrm{s})$. We also report
\begin{equation}
\rho_{\mathrm{ref}}=\frac{R}{S(\mathrm{ref})-S(P)},
\end{equation}
which expresses the gain as a fraction of the available gap to a dense or real-audio reference. Because $\rho$ is a ratio of estimated quantities, its interval can be wide when recovery is small or the reference gap is noisy.

% ========================== FIGURE 1 SPECIFICATION ==========================
% IMPORTANT: The following figure is intentionally a layout placeholder. Do not
% replace the reserved boxes with a synthetic graph. Generate the final figure
% from the committed per-prompt result artifacts.
%
% Intended final figure footprint:
%   - one figure* spanning both columns
%   - two side-by-side panels
%   - each plot area exactly about 0.49\textwidth wide and 4.80 cm high
%   - caption below, matching the current placeholder footprint
%
% PANEL (a), operating-point response of the released recovery
%   Scientific purpose:
%     Show the main duration dependence and the uncertainty beyond 10.24 s.
%   X axis:
%     requested duration in seconds: 3.84, 5.12, 7.68, 10.24, 15.36.
%   Y axis:
%     paired CLAP recovery gain R = CLAP(P+FT)-CLAP(P).
%   Main series:
%     severity-2 released 10^6-step recovery at 3.84/5.12/7.68/10.24,
%     n=192, point estimate with 95% prompt-bootstrap CI.
%   Matched extension:
%     plot the n=96 matched-subset point at 10.24 s separately (open marker)
%     and connect THAT point, not the n=192 point, to the 15.36 s n=96 result
%     with a dashed segment. This visually respects the paired D4 comparison.
%   Annotation:
%     D4 = R(15.36)-R(10.24 matched) = +0.021 [-0.023,+0.067].
%     Use wording "no clear increase" or only the CI. Do not label plateau.
%   Optional reference cue:
%     a light vertical guide at 10.24 s may mark the native duration, but it
%     must not imply that 10.24 s is a maximum.
%
% PANEL (b), symmetric training-duration intervention
%   Scientific purpose:
%     Make the specialization prediction and its falsification visually obvious.
%   X axis:
%     evaluation duration: 3.84 and 10.24 s.
%   Y axis:
%     20k-step fine-tuning gain over the SAME severity-2 P baseline.
%   Series 1:
%     shortft, trained at 3.84 s, R at both evaluation durations with 95% CIs.
%   Series 2:
%     longft, trained at 10.24 s, R at both evaluation durations with 95% CIs.
%   Annotate or bracket the two slopes:
%     J_short = +0.065 [+0.044,+0.087]
%     J_long  = +0.031 [+0.010,+0.052]
%     Delta J = J_long-J_short = -0.035 [-0.064,-0.004].
%   Visual logic:
%     specialization predicts the long-trained line should have a larger
%     long-minus-short slope. The observed ordering is the opposite.
%   Do NOT add the dense 2x2 or public dense text-FT reference to this panel.
%     They answer different questions and would obscure the matched comparison.
%   Styling:
%     simple point-lines with vertical 95% CI whiskers, no decorative shading,
%     direct labels or a small legend, readable at two-column ICASSP print size.
% ===========================================================================
\begin{figure*}[t]
\centering
\begin{minipage}[t]{0.49\textwidth}
\centering
\includegraphics[width=0.965\linewidth]{figs/fig1a_duration.pdf}
\end{minipage}\hfill
\begin{minipage}[t]{0.49\textwidth}
\centering
\includegraphics[width=0.965\linewidth]{figs/fig1b_intervention.pdf}
\end{minipage}
\caption{Duration dependence and the symmetric training-duration test. Left, the released severity~2 recovery is evaluated across duration, with the $15.36$\,s extension compared on its matched $96$-prompt subset. Right, two $20{,}000$-step fine-tunes from the same pruned checkpoint differ only in training duration and are evaluated at both $3.84$ and $10.24$\,s.}
\label{fig:duration}
\end{figure*}
